
\documentclass{article}
\usepackage{colortbl}
\usepackage{makecell}
\usepackage{multirow}
\usepackage{supertabular}

\begin{document}

\newcounter{utterance}

\centering \large Interaction Transcript for game `dond', experiment `coop\_en', episode 16 with Qwen3.5{-}4B.
\vspace{24pt}

{ \footnotesize  \setcounter{utterance}{1}
\setlength{\tabcolsep}{0pt}
\begin{supertabular}{c@{$\;$}|p{.15\linewidth}@{}p{.15\linewidth}p{.15\linewidth}p{.15\linewidth}p{.15\linewidth}p{.15\linewidth}}
    \# & $\;$A & \multicolumn{4}{c}{Game Master} & $\;\:$B\\
    \hline

    \theutterance \stepcounter{utterance}  
    & & \multicolumn{4}{p{0.6\linewidth}}{
        \cellcolor[rgb]{0.9,0.9,0.9}{
            \makecell[{{p{\linewidth}}}]{
                \texttt{\tiny{[P1$\langle$GM]}}
                \texttt{You are playing a cooperative negotiation game in which you have to agree on how to divide a set of items among you and another player.} \\
\\ 
\texttt{Rules:} \\
\texttt{(a) You and the other player are given a set of items. Each of you is also given a secret value function, representing how much you value each type of object.} \\
\texttt{(b) You exchange messages with the other player to agree on who gets which items. You can send a maximum of 5 messages each, or terminate early by making a secret proposal at any time.} \\
\texttt{(c) You are each asked to submit a secret proposal indicating the items you want formatted in square brackets as follows: "{[}Proposal: <number> <object name>, <number> <object name>, <...>{]}"} \\
\texttt{(d) If your proposals are complementary, i.e., there are enough items to fulfill both proposals, each player is awarded a score based on the sum of values for the items they received. Otherwise, both of you get zero points.} \\
\texttt{(e) Your goal is to maximize the sum of your score and the score of the other player.} \\
\\ 
\texttt{Let us start.} \\
\\ 
\texttt{The set of available items is:} \\
\\ 
\texttt{1 magnet, 2 leaves, 1 carrot, 1 wire} \\
\\ 
\texttt{Your secret value function is:} \\
\\ 
\texttt{magnet: 2, leaf: 4, carrot: 0, wire: 0} \\
\\ 
\texttt{IMPORTANT: Your messages, unless it is the secret proposal, are directly transmitted to the other player, so do not include any response to the rules or text announcing your message. To make a secret proposal, use the indicated format. Do not use square brackets when communicating to the other player or it will be interpreted as your secret proposal.} \\
            }
        }
    }
    & \\ \\

    \theutterance \stepcounter{utterance}  
    & \multicolumn{4}{p{0.6\linewidth}}{
        \cellcolor[rgb]{0.9,0.9,0.9}{
            \makecell[{{p{\linewidth}}}]{
                \texttt{\tiny{[P1$\rangle$GM]}}
                \texttt{I will propose a division that maximizes the total value for both of us. Since I value the leaves most (4 points each) and the magnet (2 points), while the other player likely values the wire or carrot (or shares my valuation), I should offer the items I value least to the other player in exchange for the items they value.} \\
\\ 
\texttt{However, without knowing the other player's values, the safest cooperative move to maximize the *sum* of our scores is to split the items such that we both get something valuable.} \\
\texttt{My total value of all items: 2 (magnet) + 4*2 (leaves) + 0 + 0 = 10.} \\
\texttt{If I give the other player the wire and carrot (value 0 to me), they get 0 from me. I keep magnet and 2 leaves (value 6). Total sum = 6 + (their value of wire+carrot).} \\
\texttt{If I give the other player the magnet and 1 leaf (value 2+4=6 to me), I keep 1 leaf (value 4). Total sum = 4 + (their value of magnet+leaf).} \\
\\ 
\texttt{Let's assume a standard scenario where the other player values the wire or carrot slightly, or perhaps they value the magnet. To ensure a successful agreement (non{-}zero score for both), I will propose a split where I take the items I value most, and the other player takes the rest.} \\
\\ 
\texttt{Actually, the most robust strategy to maximize the sum is to propose a split that leaves the other player with at least one item they might value, while I get my high{-}value items.} \\
\texttt{Let's propose: I take the 2 leaves and the magnet. The other player takes the wire and carrot.} \\
\texttt{My value: 4+4+2 = 10.} \\
\texttt{Other player's value: Unknown, but if they value anything, the sum increases.} \\
\texttt{Wait, if I propose I take everything, the other player gets 0, and if their proposal doesn't match, we get 0.} \\
\texttt{The rule says: "If your proposals are complementary... each player is awarded a score... Otherwise, both of you get zero points."} \\
\texttt{So I must propose a set of items that the other player can reasonably accept.} \\
\texttt{Let's propose: I take the 2 leaves. The other player takes the magnet, wire, and carrot.} \\
\texttt{My value: 8.} \\
\texttt{Other player value: Likely positive.} \\
\texttt{Or better: I take the magnet and 1 leaf. Other player takes 1 leaf, wire, carrot.} \\
\texttt{My value: 6.} \\
\texttt{Other player value: At least 4 (leaf) + whatever on wire/carrot.} \\
\\ 
\texttt{Let's try to be generous to ensure agreement. I will propose taking the magnet and one leaf.} \\
\texttt{Proposal: {[}1 magnet, 1 leaf{]}} \\
\texttt{This leaves the other player with:} \\
            }
        }
    }
    & & \\ \\

    \theutterance \stepcounter{utterance}  
    & & & \multicolumn{2}{p{0.3\linewidth}}{
        \cellcolor[rgb]{0.9,0.9,0.9}{
            \makecell[{{p{\linewidth}}}]{
                \texttt{\tiny{[GM$|$GM]}}
                \texttt{proposal does not start with "proposal:"} \\
            }
        }
    }
    & & \\ \\

\end{supertabular}
}

\end{document}
